\documentclass[11pt]{article}
\usepackage[margin=1in]{geometry}
\usepackage[T1]{fontenc}
\usepackage[utf8]{inputenc}
\usepackage{array}
\usepackage{longtable}
\usepackage{booktabs}
\usepackage{xcolor}
\usepackage{listings}
\usepackage{hyperref}

\hypersetup{colorlinks=true, linkcolor=blue, urlcolor=blue}

\lstdefinestyle{source}{
  basicstyle=\ttfamily\scriptsize,
  breaklines=true,
  columns=fullflexible,
  keepspaces=true,
  frame=single,
  framerule=0.2pt,
  rulecolor=\color{black!30},
  numbers=left,
  numberstyle=\tiny\color{black!45},
  xleftmargin=2em,
  framexleftmargin=1.5em
}
\lstset{style=source}

\newcommand{\src}[1]{\texttt{\detokenize{#1}}}
\newcommand{\itemhead}[4]{%
  \subsection{#1}%
  \noindent\textbf{Book source:} \src{#2}.\\
  \textbf{Formal source:} \src{#3}.\\
  \textbf{Formal item:} \src{#4}.%
}

\title{Moise Sections 3--4: Detailed Formalization Correspondence}
\author{GeoTop formalization notes}
\date{Snapshot: 2026-06-17}

\begin{document}
\maketitle

\section{How to Read This File}

This report expands \src{SECTIONS_3_4_FORMALIZATION_MAP.md}.  For each book
item in Sections 3--4 it gives:

\begin{itemize}
  \item the exact book source range from \src{geotop.tex}, including the book
    statement and proof text for that item;
  \item the exact current Isabelle declaration range, with file and line
    number;
  \item a short explanation of how the Isabelle proof was organized and how it
    corresponds to Moise's proof.
\end{itemize}

The ``exact'' excerpts are inserted with \verb|\lstinputlisting| from the live
files in \src{tst/}.  Compile this file from \src{/project/tst} so the relative
paths resolve.

\section{Section 3: The Schoenflies Theorem for Polygons}

\itemhead{Theorem 3.1}{geotop.tex:724--733}{dev34_prefix_base/GeoTop_3_4_Prefix_Base.thy:551}{Theorem_GT_3_1}

\paragraph{Book statement and proof.}
\lstinputlisting[firstline=724,lastline=733]{geotop.tex}

\paragraph{Formal statement.}
\lstinputlisting[firstline=548,lastline=563]{dev34_prefix_base/GeoTop_3_4_Prefix_Base.thy}

\paragraph{Isabelle proof and correspondence.}
The Isabelle proof follows the barycentric-coordinate proof, but packages the
hard affine extension in existing convex/simplex infrastructure.  The theorem is
stated for an arbitrary Euclidean space type, with vertex sets \src{V}, \src{W}
and a bijection \src{\<phi>}.  The proof first exposes barycentric coordinates
for points of \src{\<sigma>}, then uses
\src{geotop_simplex_vertex_bijection_affine_extension} to obtain the affine
homeomorphism matching the prescribed vertices.  It bridges the book's
``simplicial homeomorphism'' to \src{top1_homeomorphism_on} and
\src{geotop_simplicial_on}.

\itemhead{Theorem 3.2}{geotop.tex:735--750}{dev34_prefix_base/GeoTop_3_4_Prefix_Base.thy:721}{Theorem_GT_3_2}

\paragraph{Book statement and proof.}
\lstinputlisting[firstline=735,lastline=750]{geotop.tex}

\paragraph{Formal statement.}
\lstinputlisting[firstline=718,lastline=729]{dev34_prefix_base/GeoTop_3_4_Prefix_Base.thy}

\paragraph{Isabelle proof and correspondence.}
Moise translates the simplex to the origin and extends the vertex map linearly
to the ambient space.  Isabelle reuses Theorem 3.1 for the restricted simplex
homeomorphism and then applies
\src{geotop_simplex_vertex_bijection_affine_extension} to get the global
Euclidean homeomorphism \src{g}.  The formal statement records the same
ambient-space conclusion as a homeomorphism of \src{UNIV}.

\itemhead{Definition: Free 2-Simplex}{geotop.tex:752--760}{dev34_prefix_mid/GeoTop_3_4_Prefix_Mid.thy:5808,5817}{geotop_free_2_simplex; geotop_boundary_free_2_simplex}

\paragraph{Book definition.}
\lstinputlisting[firstline=752,lastline=760]{geotop.tex}

\paragraph{Formal definitions.}
\lstinputlisting[firstline=5804,lastline=5828]{dev34_prefix_mid/GeoTop_3_4_Prefix_Mid.thy}

\paragraph{Isabelle proof and correspondence.}
The book says a free triangle is one whose intersection with the polygon
frontier consists of one or two edges.  The formal definition also permits the
empty-contact case because later fold-normalization needs to distinguish a
chosen free 2-simplex from one whose actual boundary contact is nonempty.
\src{geotop_boundary_free_2_simplex} is the nonempty boundary-contact version.
The surrounding lemmas prove transfer, cardinality, and contact-case facts that
let the induction in Theorem 3.3 recover the book's intended free-boundary
witness.

\itemhead{Theorem 3.3}{geotop.tex:762--780}{dev34_prefix_mid/GeoTop_3_4_Prefix_Mid.thy:19301}{Theorem_GT_3_3}

\paragraph{Book statement and proof.}
\lstinputlisting[firstline=762,lastline=780]{geotop.tex}

\paragraph{Formal statement.}
\lstinputlisting[firstline=19298,lastline=19318]{dev34_prefix_mid/GeoTop_3_4_Prefix_Mid.thy}

\paragraph{Isabelle proof and correspondence.}
The final theorem wrapper is short: it invokes the stronger count theorem
\src{geotop_polygon_disk_free_2simplex_count_ge2_prefix}, exactly matching
Moise's proof strategy that proves at least two free 2-simplexes.  The long
formal work is in the induction infrastructure: existence of two boundary
2-simplexes, the nonfree boundary-triangle split along the chord, decomposition
into two smaller polygonal disks, and transfer of side witnesses back to the
parent disk.  The last extraction step is
\src{geotop_polygon_disk_boundary_free_witness_avoids_empty_contact_prefix}.

\itemhead{Theorem 3.4}{geotop.tex:782--799}{dev34_prefix_mid/GeoTop_3_4_Prefix_Mid.thy:58990}{Theorem_GT_3_4}

\paragraph{Book statement and proof.}
\lstinputlisting[firstline=782,lastline=799]{geotop.tex}

\paragraph{Formal statement.}
\lstinputlisting[firstline=58987,lastline=59004]{dev34_prefix_mid/GeoTop_3_4_Prefix_Mid.thy}

\paragraph{Isabelle proof and correspondence.}
The formal proof obtains a triangulation of the closed polygonal disk via
\src{Theorem_GT_2_2}, then runs the free-triangle deletion induction as a
supported fold-normalization theorem.  The two Figure 3.3 cases are formalized
by \src{geotop_figure33_one_boundary_named_supported_fold_prefix} and
\src{geotop_figure33_two_boundary_named_supported_inverse_fold_prefix}; the
combined induction package is
\src{geotop_polygon_disk_free_triangle_fold_normalization_supported_prefix}.
This is the direct formal replacement for Moise's ``remove a free
2-simplex by a homeomorphism'' step.

\itemhead{Theorem 3.5}{geotop.tex:801--820}{dev34_prefix_mid/GeoTop_3_4_Prefix_Mid.thy:59533}{Theorem_GT_3_5}

\paragraph{Book statement and proof.}
\lstinputlisting[firstline=801,lastline=820]{geotop.tex}

\paragraph{Formal statement.}
\lstinputlisting[firstline=59530,lastline=59543]{dev34_prefix_mid/GeoTop_3_4_Prefix_Mid.thy}

\paragraph{Isabelle proof and correspondence.}
The Isabelle proof follows the composition proof in the book.  It applies
\src{Theorem_GT_3_4} to \src{J} and \src{J'} to send both polygons to simplex
frontiers, applies the simplex homeomorphism machinery from Theorem 3.2 to
match those frontiers, then composes the three plane homeomorphisms using the
project's homeomorphism-composition and inverse lemmas.

\itemhead{Theorem 3.6}{geotop.tex:821--822}{dev34_prefix_mid/GeoTop_3_4_Prefix_Mid.thy:59717}{Theorem_GT_3_6}

\paragraph{Book statement and proof.}
\lstinputlisting[firstline=821,lastline=822]{geotop.tex}

\paragraph{Formal statement.}
\lstinputlisting[firstline=59715,lastline=59726]{dev34_prefix_mid/GeoTop_3_4_Prefix_Mid.thy}

\paragraph{Isabelle proof and correspondence.}
Moise says simply ``By Theorem 4.''  Isabelle expands that into the standard
preimage argument: from Theorem 3.4 choose a homeomorphism carrying \src{J} to
the frontier of a 2-simplex \src{\<sigma>}; a simplex is a 2-cell, and the
preimage of that 2-cell under the homeomorphism is the required disk
\src{D}.  Frontier preservation under plane homeomorphisms supplies
\src{J = Fr D}.

\itemhead{Theorem 3.7}{geotop.tex:824--826}{dev34_prefix_mid/GeoTop_3_4_Prefix_Mid.thy:59831}{Theorem_GT_3_7}

\paragraph{Book statement and proof.}
\lstinputlisting[firstline=824,lastline=826]{geotop.tex}

\paragraph{Formal statement.}
\lstinputlisting[firstline=59824,lastline=59860]{dev34_prefix_mid/GeoTop_3_4_Prefix_Mid.thy}

\paragraph{Isabelle proof and correspondence.}
The book's one-line proof says to choose the deletion homeomorphisms with
support inside \src{U}.  Isabelle makes this explicit by using the same
fold-normalization package as Theorem 3.4, but in its support-parametric form:
\src{geotop_polygon_disk_free_triangle_fold_normalization_supported_prefix}.
The formal conclusion adds \src{\<forall>P\<in>UNIV - U. h P = P}, exactly encoding
the fixed-outside-\src{U} requirement.

\section{Section 4: The Jordan Curve Theorem}

\itemhead{Theorem 4.1}{geotop.tex:865--867}{dev34_prefix_mid/GeoTop_3_4_Prefix_Mid.thy:59930}{Theorem_GT_4_1}

\paragraph{Book statement and proof.}
\lstinputlisting[firstline=865,lastline=867]{geotop.tex}

\paragraph{Formal statement.}
\lstinputlisting[firstline=59930,lastline=59942]{dev34_prefix_mid/GeoTop_3_4_Prefix_Mid.thy}

\paragraph{Isabelle proof and correspondence.}
The formal proof sets \src{V} to the component of \src{U} containing \src{P}
and \src{W = U - V}.  It proves the component is open by
\src{geotop_component_at_open_in_euclidean}, proves the complement is a union
of open components, and uses component disjointness to place \src{Q} in
\src{W}.  This is Moise's proof with the component-open fact made explicit.

\itemhead{Theorem 4.2}{geotop.tex:869--884}{dev34_prefix_mid/GeoTop_3_4_Prefix_Mid.thy:67224}{Theorem_GT_4_2}

\paragraph{Book statement and proof.}
\lstinputlisting[firstline=869,lastline=884]{geotop.tex}

\paragraph{Formal statement.}
\lstinputlisting[firstline=67219,lastline=67255]{dev34_prefix_mid/GeoTop_3_4_Prefix_Mid.thy}

\paragraph{Isabelle proof and correspondence.}
The formal theorem includes an explicit cyclic-order assumption
\src{geotop_polygon_cyclic_order J P Q R S}.  The proof delegates the geometric
contradiction argument to
\src{geotop_polygon_arc_opposite_boundary_decomposition_prefix}, which bundles
the reduction to the rectangle, the ``if Q' and S' were in the same component''
broken-line construction, and the contradiction with the earlier crossing
theorem.  The wrapper then reorders the conjuncts into the theorem statement.

\itemhead{Theorem 4.3}{geotop.tex:886--929}{dev34_prefix/GeoTop_3_4_Prefix.thy:7}{Theorem_GT_4_3}

\paragraph{Book statement and proof, including Lemmas 4.3.1 and 4.3.2.}
\lstinputlisting[firstline=886,lastline=929]{geotop.tex}

\paragraph{Formal statement.}
\lstinputlisting[firstline=5,lastline=34]{dev34_prefix/GeoTop_3_4_Prefix.thy}

\paragraph{Isabelle proof and correspondence.}
The book gives a polygonal-cell proof with two internal lemmas.  The current
formal proof does not name those sublemmas separately; it proves the same
separation theorem by bridging \src{geotop_is_n_sphere} to HOL-Analysis
homeomorphism of the standard circle and applying
\src{Jordan_Brouwer_separation}.  Thus Lemma 4.3.1 and Lemma 4.3.2 are
absorbed into the library theorem rather than replayed as separate Isabelle
lemmas.

\paragraph{Lemma 4.3.1 row.}
The exact book statement and proof are \src{geotop.tex:898--902}, included in
the listing above.  There is no standalone Isabelle declaration for this
sublemma; its separation role is covered by \src{Theorem_GT_4_3} at
\src{dev34_prefix/GeoTop_3_4_Prefix.thy:7}.

\paragraph{Lemma 4.3.2 row.}
The exact book statement and proof are \src{geotop.tex:904--929}, included in
the listing above.  As with Lemma 4.3.1, the current formalization proves the
parent theorem via \src{Jordan_Brouwer_separation}; no separate Isabelle lemma
is introduced for the bounded-component construction.

\itemhead{Theorem 4.4}{geotop.tex:931--974}{dev34_prefix/GeoTop_3_4_Prefix.thy:32828}{Theorem_GT_4_4}

\paragraph{Book statement and proof.}
\lstinputlisting[firstline=931,lastline=974]{geotop.tex}

\paragraph{Formal statement.}
\lstinputlisting[firstline=32824,lastline=32857]{dev34_prefix/GeoTop_3_4_Prefix.thy}

\paragraph{Isabelle proof and correspondence.}
The formal theorem is stated for a polygon \src{J} with named boundary points
\src{P,Q,R,S}, cyclic order, two disjoint arcs \src{A1,A2}, and the endpoint
conditions \src{A1 \<inter> J = {P}} and \src{A2 \<inter> J = {R}}.  The proof is
the one-line wrapper around
\src{geotop_polygon_two_disjoint_endpoint_arcs_brick_component_transfer_prefix}.
That auxiliary lemma is the formal version of the brick/regular-neighborhood
construction: choose a sufficiently fine carrier around \src{A1}, take the
frontier component containing \src{P}, extract the broken line crossing from
the lower to upper side, and transfer the resulting component-frontier fact to
\src{Q} and \src{S}.

\itemhead{Definition: Brick-Decomposition}{geotop.tex:943--956}{dev34_prefix/GeoTop_3_4_Prefix.thy:41}{geotop_brick_decomposition}

\paragraph{Book definition.}
\lstinputlisting[firstline=943,lastline=956]{geotop.tex}

\paragraph{Formal definition.}
\lstinputlisting[firstline=36,lastline=50]{dev34_prefix/GeoTop_3_4_Prefix.thy}

\paragraph{Isabelle proof and correspondence.}
The definition records exactly the three book requirements: each brick is a
polyhedral disk, the union is all of the plane, pairwise nonempty intersections
are broken lines in both frontiers, and every point has a neighborhood meeting
at most three bricks.  The completed Theorem 4.4 proof uses a specialized
transfer lemma rather than developing a broad reusable theory of all brick
decompositions.

\itemhead{Theorem 4.5}{geotop.tex:976--994}{dev34_facts/GeoTop_3_4_Facts.thy:11}{Theorem_GT_4_5}

\paragraph{Book statement and proof.}
\lstinputlisting[firstline=976,lastline=994]{geotop.tex}

\paragraph{Formal statement.}
\lstinputlisting[firstline=9,lastline=18]{dev34_facts/GeoTop_3_4_Facts.thy}

\paragraph{Isabelle proof and correspondence.}
Moise derives this from Theorem 4.4 by a bounded-component contradiction.  The
formal proof uses the already available general result
\src{top1_connected_complement_of_geotop_arc}, specialized to dimension two.
This is stronger than the book route but proves the same theorem with the same
geometric content: an arc complement in the plane is connected.

\itemhead{Theorem 4.6}{geotop.tex:996--998}{dev34_facts/GeoTop_3_4_Facts.thy:22}{Theorem_GT_4_6}

\paragraph{Book statement and proof.}
\lstinputlisting[firstline=996,lastline=998]{geotop.tex}

\paragraph{Formal statement.}
\lstinputlisting[firstline=20,lastline=29]{dev34_facts/GeoTop_3_4_Facts.thy}

\paragraph{Isabelle proof and correspondence.}
The book argues that a proper frontier subset would lie in an arc and separate
the plane, contradicting Theorem 4.5.  Isabelle proves the same frontier
conclusion by using \src{Jordan_Brouwer_frontier} after bridging a geotop
1-sphere to the HOL standard sphere.  The component hypothesis is expressed via
\src{geotop_component_at}.

\itemhead{Theorem 4.7}{geotop.tex:1002--1016}{dev34_facts/GeoTop_3_4_Facts.thy:265}{Theorem_GT_4_7}

\paragraph{Book statement and proof.}
\lstinputlisting[firstline=1002,lastline=1016]{geotop.tex}

\paragraph{Formal statement.}
\lstinputlisting[firstline=263,lastline=279]{dev34_facts/GeoTop_3_4_Facts.thy}

\paragraph{Isabelle proof and correspondence.}
Moise proves uniqueness of the bounded component using the same Figure 4.2/4.9
configuration and Theorem 4.5.  The Isabelle proof splits the argument into
existence of a bounded component and uniqueness.  Existence comes from
\src{geotop_polygon_interior_is_bounded_component}; uniqueness uses the same
``a second bounded component would have frontier in an arc'' contradiction,
with the arc-separation obstruction supplied by Theorem 4.5 and Jordan-Brouwer
component facts.

\itemhead{Theorem 4.8 and Lemmas 4.8.1--4.8.5}{geotop.tex:1020--1051}{dev34/GeoTop_3_4.thy:28666}{Theorem_GT_4_8}

\paragraph{Book statement and proof.}
\lstinputlisting[firstline=1020,lastline=1051]{geotop.tex}

\paragraph{Formal statement.}
\lstinputlisting[firstline=28663,lastline=28682]{dev34/GeoTop_3_4.thy}

\paragraph{Isabelle proof and correspondence.}
The formal proof follows the book's five-lemma outline inside the theorem body.
Local fact \src{hL1} proves every vertex lies in an edge using the no-open-
singleton property of a 2-manifold.  Local \src{hL2} proves at least one
incident 2-simplex for every incident edge.  Local \src{hL3} proves at least
two incident 2-simplexes using the small-semicircle/no-arc-separates-plane
package.  Local \src{hL4} proves at most two using the three-incident-simplex
Jordan contradiction.  Local \src{hL5} proves link connectedness.  The final
step uses the Figure 4.10 cone/link model lemmas to identify each star with a
combinatorial 2-cell.

\paragraph{Lemma 4.8.1 row.}
Book source: \src{geotop.tex:1024--1025}.  Formal status:
\src{geotop_complex_simplex_vertex_incident_edge} at
\src{dev34_facts/GeoTop_3_4_Facts.thy:535}, and the local theorem step
\src{hL1} at \src{dev34/GeoTop_3_4.thy:28686}.

\paragraph{Lemma 4.8.2 row.}
Book source: \src{geotop.tex:1027--1029}.  Formal status:
\src{geotop_incident_edge_face_count_ge_1_link_edge_witness} at
\src{dev34_facts/GeoTop_3_4_Facts.thy:4892}, and local steps \src{hL2} and
\src{hL2_count} at \src{dev34/GeoTop_3_4.thy:28715} and
\src{dev34/GeoTop_3_4.thy:28739}.

\paragraph{Lemma 4.8.3 row.}
Book source: \src{geotop.tex:1031--1032}.  Formal status:
\src{geotop_unique_incident_2simplex_small_semicircle_separates_chart_dev34}
at \src{dev34/GeoTop_3_4.thy:24605}, with the local conclusion \src{hL3} at
\src{dev34/GeoTop_3_4.thy:28770}.

\paragraph{Lemma 4.8.4 row.}
Book source: \src{geotop.tex:1034--1035}.  Formal status:
\src{geotop_three_incident_2simplex_small_circle_not_separates_chart_dev34}
at \src{dev34/GeoTop_3_4.thy:27813},
\src{geotop_three_incident_2simplex_sphere_not_separates_chart_dev34} at
\src{dev34/GeoTop_3_4.thy:27888}, and
\src{geotop_boundary_2cell_chart_three_incident_2simplex_contradiction_dev34}
at \src{dev34/GeoTop_3_4.thy:28588}.  The local conclusion is \src{hL4} at
\src{dev34/GeoTop_3_4.thy:28850}.

\paragraph{Lemma 4.8.5 row.}
Book source: \src{geotop.tex:1039--1040}.  Formal status: the link
connectedness conclusion is the local fact \src{hL5} at
\src{dev34/GeoTop_3_4.thy:29543}; it is not split out as a global theorem.

\paragraph{Figure 4.10 model row.}
Book source: \src{geotop.tex:1042--1051}.  Formal status:
\src{geotop_fig410_boundary_subdivision_model_from_finite_linear_link_polygon_dev34}
at \src{dev34/GeoTop_3_4.thy:7616},
\src{geotop_standard_boundary_cone_model_from_finite_linear_link_polygon_dev34}
at \src{dev34/GeoTop_3_4.thy:12048}, and
\src{geotop_standard_boundary_cone_model_from_finite_linear_link_polygon_with_interior_dev34}
at \src{dev34/GeoTop_3_4.thy:12142}.

\itemhead{Theorem 4.9 and Combinatorial Boundary Definition}{geotop.tex:1052--1056}{dev34/GeoTop_3_4.thy:31656}{Theorem_GT_4_9}

\paragraph{Book statement and proof.}
\lstinputlisting[firstline=1052,lastline=1056]{geotop.tex}

\paragraph{Formal statement.}
\lstinputlisting[firstline=31652,lastline=31669]{dev34/GeoTop_3_4.thy}

\paragraph{Isabelle proof and correspondence.}
The proof reuses the 4.8 link/star analysis with the boundary case: each edge
has one or two incident 2-simplexes, and the link of a vertex is either a
broken line or a polygon.  The formal boundary assertion is stated directly as
the union of edges with exactly one incident 2-simplex.  The key named bridge
is \src{geotop_manifold_boundary_eq_one_incident_edges_dev34}, preceded by
\src{geotop_edge_face_count_one_or_two_in_manifold_with_boundary_dev34}.

\itemhead{Theorem 4.10}{geotop.tex:1058--1062}{dev34/GeoTop_3_4.thy:32151}{Theorem_GT_4_10}

\paragraph{Book statement and proof.}
\lstinputlisting[firstline=1058,lastline=1062]{geotop.tex}

\paragraph{Formal statement.}
\lstinputlisting[firstline=32148,lastline=32163]{dev34/GeoTop_3_4.thy}

\paragraph{Isabelle proof and correspondence.}
The formal proof follows Moise's two inclusions.  For
\src{Bd M \<subseteq> Fr M}, a point of \src{M - Fr M} has an open neighborhood in
\src{M} and is therefore manifold-interior.  For \src{Fr M \<subseteq> Bd M}, closedness
puts frontier points in \src{M}; if such a point were manifold-interior, its
local chart would give a plane-homeomorphic open neighborhood, and
\src{Theorem_GT_4_invariance_of_domain} would make that neighborhood open in
\src{R^2}, contradicting membership in the frontier.

\paragraph{Invariance-of-domain support.}
\lstinputlisting[firstline=6715,lastline=6728]{gb/GeoTopBase.thy}

\section{Main Auxiliary Isabelle Lemmas}

\begin{longtable}{>{\raggedright\arraybackslash}p{0.28\textwidth}
                  >{\raggedright\arraybackslash}p{0.25\textwidth}
                  >{\raggedright\arraybackslash}p{0.39\textwidth}}
\toprule
\textbf{Isabelle item} & \textbf{Location} & \textbf{Book correspondence / role}\\
\midrule
\endhead
\src{geotop_polygon_disk_two_boundary_2simplexes_prefix} &
\src{dev34_prefix_mid/GeoTop_3_4_Prefix_Mid.thy:5493} &
Supports Theorem 3.3: supplies the book's claim that at least two 2-simplexes
touch the polygon frontier.\\
\src{geotop_polygon_disk_nonfree_boundary_triangle_split_free_count_prefix} &
\src{dev34_prefix_mid/GeoTop_3_4_Prefix_Mid.thy:12903} &
Formalizes the Figure 3.2 split when the chosen boundary triangle is not free.\\
\src{geotop_polygon_disk_nonfree_boundary_triangle_decomposition_free_count_prefix} &
\src{dev34_prefix_mid/GeoTop_3_4_Prefix_Mid.thy:13771} &
Runs the two-side induction count used in Moise's stronger ``two free
2-simplexes'' argument.\\
\src{geotop_polygon_disk_boundary_free_witness_avoids_empty_contact_prefix} &
\src{dev34_prefix_mid/GeoTop_3_4_Prefix_Mid.thy:34363} &
Extracts a genuine boundary-free witness distinct from the chosen bad triangle.\\
\src{geotop_figure33_one_boundary_named_supported_fold_prefix} &
\src{dev34_prefix_mid/GeoTop_3_4_Prefix_Mid.thy:45304} &
Formal Figure 3.3 Case 1 fold, with support control.\\
\src{geotop_figure33_two_boundary_named_supported_inverse_fold_prefix} &
\src{dev34_prefix_mid/GeoTop_3_4_Prefix_Mid.thy:55144} &
Formal Figure 3.3 Case 2 inverse fold, with support control.\\
\src{geotop_free_triangle_one_boundary_edge_supported_fold_prefix} &
\src{dev34_prefix_mid/GeoTop_3_4_Prefix_Mid.thy:58113} &
Packages the one-boundary-edge deletion case for Theorems 3.4 and 3.7.\\
\src{geotop_free_triangle_two_boundary_edges_supported_inverse_fold_prefix} &
\src{dev34_prefix_mid/GeoTop_3_4_Prefix_Mid.thy:58242} &
Packages the two-boundary-edge deletion case for Theorems 3.4 and 3.7.\\
\src{geotop_polygon_disk_free_triangle_fold_normalization_supported_prefix} &
\src{dev34_prefix_mid/GeoTop_3_4_Prefix_Mid.thy:58843} &
The full supported induction that removes free triangles until the polygon is
carried to a simplex frontier.\\
\src{geotop_polygon_arc_opposite_boundary_decomposition_prefix} &
\src{dev34_prefix_mid/GeoTop_3_4_Prefix_Mid.thy:67081} &
Formalizes the core contradiction in Theorem 4.2.\\
\src{geotop_polygon_two_disjoint_endpoint_arcs_brick_component_transfer_prefix} &
\src{dev34_prefix/GeoTop_3_4_Prefix.thy:30025} &
Formalizes the brick/regular-neighborhood component-transfer proof of
Theorem 4.4.\\
\src{geotop_complex_simplex_vertex_incident_edge} &
\src{dev34_facts/GeoTop_3_4_Facts.thy:535} &
Standalone support for Lemma 4.8.1: vertices are incident to edges in the
2-manifold setting.\\
\src{geotop_incident_edge_face_count_ge_1_link_edge_witness} &
\src{dev34_facts/GeoTop_3_4_Facts.thy:4892} &
Support for Lemma 4.8.2: incident edges have a 2-simplex witness.\\
\src{geotop_unique_incident_2simplex_small_semicircle_separates_chart_dev34} &
\src{dev34/GeoTop_3_4.thy:24605} &
The small-semicircle contradiction behind Lemma 4.8.3.\\
\src{geotop_three_incident_2simplex_small_circle_not_separates_chart_dev34} &
\src{dev34/GeoTop_3_4.thy:27813} &
Initial local contradiction for three 2-simplexes incident to one edge.\\
\src{geotop_three_incident_2simplex_sphere_not_separates_chart_dev34} &
\src{dev34/GeoTop_3_4.thy:27888} &
Jordan/sphere version of the three-incident-simplex obstruction.\\
\src{geotop_boundary_2cell_chart_three_incident_2simplex_contradiction_dev34} &
\src{dev34/GeoTop_3_4.thy:28588} &
Final contradiction package used in Lemma 4.8.4.\\
\src{geotop_fig410_boundary_subdivision_model_from_finite_linear_link_polygon_dev34} &
\src{dev34/GeoTop_3_4.thy:7616} &
Builds the Figure 4.10 boundary subdivision model.\\
\src{geotop_standard_boundary_cone_model_from_finite_linear_link_polygon_dev34} &
\src{dev34/GeoTop_3_4.thy:12048} &
Identifies a star with the standard cone over a polygonal link.\\
\src{geotop_standard_boundary_cone_model_from_finite_linear_link_polygon_with_interior_dev34} &
\src{dev34/GeoTop_3_4.thy:12142} &
Interior-aware version used to conclude the star is a combinatorial 2-cell.\\
\src{geotop_edge_face_count_one_or_two_in_manifold_with_boundary_dev34} &
\src{dev34/GeoTop_3_4.thy:30251} &
Boundary-manifold edge count for Theorem 4.9.\\
\src{geotop_manifold_boundary_eq_one_incident_edges_dev34} &
\src{dev34/GeoTop_3_4.thy:31626} &
Identifies the topological manifold boundary with the union of one-incident
edges, matching Moise's \(\partial K\) paragraph.\\
\src{Theorem_GT_4_invariance_of_domain} &
\src{gb/GeoTopBase.thy:6722} &
The Brouwer invariance-of-domain theorem used in the second inclusion of
Theorem 4.10.\\
\bottomrule
\end{longtable}

\section{Summary}

Sections 3--4 are now represented by named Isabelle theorems
\src{Theorem_GT_3_1} through \src{Theorem_GT_3_7} and
\src{Theorem_GT_4_1} through \src{Theorem_GT_4_10}.  The proof style is not
uniformly literal: several book picture arguments are replaced by packaged
Isabelle lemmas, and Theorems 4.3, 4.5, and 4.6 use HOL-Analysis/Jordan-Brouwer
bridges where that is cleaner than replaying Moise's full planar construction.
The correspondences above record where those replacements occur and which
auxiliary lemmas carry the geometric work.

\end{document}
